% Background Validity and the B₀→0 Argument
% Central theoretical document for the torsion-Gertsenshtein sweep programme.
% Extends Hwang & Noh (2310.04150) to arbitrary PGT+EM theories.

\section{Background Validity for Gertsenshtein-Type Linearisation}
\label{sec:background-validity}

\subsection{Problem Statement}
\label{sec:bv-problem}

The Gertsenshtein effect is computed by linearising about a background
consisting of flat Minkowski spacetime, a uniform magnetic field~$B_0$,
and zero torsion:
\begin{equation}
  g_{\mu\nu} = \eta_{\mu\nu} + \epsilon\, h_{\mu\nu}, \qquad
  A_\mu = \bar{A}_\mu + \epsilon\, a_\mu, \qquad
  T^\lambda{}_{\mu\nu} = 0 + \epsilon\, t^\lambda{}_{\mu\nu}\,.
  \label{eq:bv-background}
\end{equation}
This background is \textbf{not} an exact solution of the full nonlinear
field equations: the magnetic field carries stress-energy
$T_{\mu\nu}^{(\mathrm{EM})} \sim B_0^2/2$ that sources curvature, and
non-minimal couplings can source background torsion from~$B_0$.

Two questions arise:
\begin{enumerate}
  \item \textbf{Backreaction:} How large are the errors from using an
    inconsistent background?
  \item \textbf{Completeness:} Which Lagrangian terms contribute at
    $O(\epsilon^2)$ when expanding around $\bar{F} \neq 0$?
\end{enumerate}

For Einstein--Maxwell, Hwang \& Noh~\cite{hwangnoh2023graviton}
(\S\,III.A, eqs.~22--31) show that the backreaction parameter
$(L/\lambda_B)^2 \sim P/(2\pi^2)$ is proportional to the conversion
probability~$P$ itself, hence negligible whenever $P \ll 1$.  However,
their analysis is specific to Einstein--Maxwell and does not address
torsion, non-minimal couplings, or higher-order EFT operators.  This
document provides a \textbf{theory-independent} generalisation.


%% ================================================================
\subsection{EFT Structure of PGT+EM}
\label{sec:bv-eft}

\subsubsection{Mass Dimension as the Organising Principle}
\label{sec:bv-mass-dim}

The most general PGT+EM Lagrangian consistent with local Lorentz and
$U(1)$ gauge symmetry is organised by the \textbf{mass dimension} of each
operator in 4D natural units ($\hbar = c = 1$, Lagrangian density~$[M^4]$).
The fundamental field strengths have dimensions:
\begin{equation}
  [\widetilde{R}_{\mu\nu\rho\sigma}] = [M^2], \qquad
  [T^\lambda{}_{\mu\nu}] = [M^1], \qquad
  [F_{\mu\nu}] = [M^2], \qquad
  [\nabla T] = [M^2]\,.
  \label{eq:bv-dimensions}
\end{equation}
where torsion has $[M^1]$ because it is an antisymmetric part of the
connection $\Gamma \sim [M^1]$.

\paragraph{Dimension-4 operators (unsuppressed).}
All operators whose contractions sum to mass dimension~4 in the Lagrangian
density have \textbf{dimensionless} coupling constants and are not
suppressed by any UV scale.  These include:
\begin{itemize}
  \item \textbf{Quadratic field-strength:} $\widetilde{R}^2$, $T^2$,
    $F^2$, $\widetilde{R} \times F$, $T \times F$, $\nabla T \times F$,
    $\widetilde{R} \times \nabla T$, $(\nabla T)^2$
    --- the standard PGT+EM action (see
    \texttt{research/general\_quadratic\_lagrangian.tex}, \S\S\,2--3).
  \item \textbf{``Cubic'' field-count, dimension-4:}
    $T^2 \times F$ (3 parity-even + 15 parity-odd contractions),
    $T^2 \times \widetilde{R}$, $T^4$
    --- these have \textbf{the same mass dimension} as the quadratic sector
    because $T$ has $[M^1]$, not $[M^2]$.  They are equally unsuppressed
    and should be included in the general theory enumeration.
\end{itemize}

\paragraph{Dimension-6 operators (suppressed by $1/\Lambda^2$).}
Operators with Lagrangian dimension~$[M^6]$ require coupling constants
with $[M^{-2}]$, parametrised as $\lambda/\Lambda^2$ where $\lambda$
is $O(1)$ and $\Lambda$ is the UV cutoff:
\begin{itemize}
  \item $\widetilde{R} \times F^2$ (4~contractions),
    $\widetilde{R}^2 \times F$ (8),
    $\nabla T \times F^2$ (5), plus parity-odd counterparts.
\end{itemize}

\paragraph{Dimension-8+ operators (doubly suppressed).}
$F^4$ (Euler--Heisenberg, coupling $\sim \alpha^2/m_e^4$),
higher-curvature terms, etc.  Always negligible for $B_0 \ll B_\mathrm{Schwinger}
\approx 4.4 \times 10^{13}\,$G.

\subsubsection{The UV Cutoff $\Lambda_4$}
\label{sec:bv-cutoff}

Aoki~\cite{aoki2020ghostfree} identifies the PGT cutoff scale:
\begin{equation}
  \Lambda_4 = (M_\mathrm{Pl}\, m^3)^{1/4}\,,
  \label{eq:bv-cutoff}
\end{equation}
where $m$ is the torsion mass (from $\alpha$ parameters).  For
$m \sim 1\,$GeV: $\Lambda_4 \sim 30\,$TeV.  Dimension-6 corrections
scale as $(E/\Lambda_4)^2 \lesssim 10^{-9}$ for gravitational-wave
physics --- entirely negligible.

\subsubsection{``How Do We Know When to Stop?''}
\label{sec:bv-truncation}

The EFT expansion answers this: enumerate all operators at mass
dimension~$\leq 4$ (a finite, calculable set).  Operators at dimension~6
are suppressed by $1/\Lambda^2$ with a \emph{known} cutoff scale.
Higher dimensions are progressively more suppressed.  This replaces the
ambiguous ``quadratic/cubic'' field-count classification with a
principled, convergent hierarchy.

The research document
(\texttt{research/general\_quadratic\_lagrangian.tex}) enumerates
$\sim 35$ quadratic terms and $\sim 18$ unsuppressed (dimension-4)
``cubic'' terms ($T^2 F$, $T^2 \widetilde{R}$, $T^4$).  An additional
$\sim 65$ dimension-6 cubic terms are suppressed by $1/\Lambda^2$ and
constitute a future extension.


%% ================================================================
\subsection{The $B_0 \to 0$ Argument}
\label{sec:bv-b0-argument}

\subsubsection{Statement}

\begin{quote}
\textbf{Theorem (informal).}\;
For any PGT+EM theory with Lagrangian of mass dimension~$\leq 4$,
linearised around flat Minkowski with uniform $B_0$ and $\bar{T} = 0$,
the conversion coefficient
$C_0 \equiv \lim_{B_0 \to 0} P(B_0)/B_0^2$
is (i)~finite and well-defined, (ii)~computed on an exactly
self-consistent background, and (iii)~free from all backreaction
and higher-order EFT corrections.
\end{quote}

\subsubsection{Proof Structure}

\paragraph{(i) Graviton--photon conversion requires $B_0$.}
The mixing coefficient $g$ in the equations of motion vanishes at
$B_0 = 0$.  Physical reason: conversion of a helicity-2 graviton to a
helicity-1 photon requires a background helicity-1 field to absorb the
angular momentum difference.  In the Lagrangian, every graviton--photon
coupling term involves at least one background factor~$\bar{F} \propto B_0$.
For the non-minimal coupling $\delta_1 \widetilde{R}_{[\mu\nu]} F^{\mu\nu}$:
the torsion-mediated chain $h \to t \to a$ uses one~$\bar{F}$ leg in the
$h \to t$ step (source $\sim \delta_1 \bar{F} \cdot \partial^2 h$);
the $t \to a$ step proceeds at perturbation level (no $\bar{F}$ needed).
Total: one power of~$B_0$.

Verified numerically: issues~\#199 and \#200 confirm zero conversion at
$B_0 = 0$ for both minimal and non-minimal PGT theories.

\paragraph{(ii) $P$ is analytic in $B_0$.}
The linearised equations of motion have polynomial dependence on $B_0$
(it enters as a coefficient in the coupling and mass terms).  The
evolution operator $M(t; B_0)$ is therefore analytic in $B_0$.  Since
$M(t; 0) = 0$ (no mixing), $P = |M|^2$ has a Taylor expansion
$P = c_2\, B_0^2 + c_4\, B_0^4 + \cdots$ with $c_2$ finite.

\paragraph{(iii) The $B_0 = 0$ background is exactly self-consistent.}
At $B_0 = 0$: $\bar{A} = 0$, $\bar{T} = 0$, $g = \eta$.  This
configuration satisfies the full nonlinear field equations of
\emph{any} PGT+EM theory trivially (all field strengths and their
stress-energy vanish).  There is no metric backreaction, no torsion
sourcing, no self-consistency error.

\paragraph{(iv) $C_0 = c_2$ is computed on the exact background.}
The coefficient $c_2 = \lim_{B_0 \to 0} P/B_0^2$ is obtained by
evaluating the linearised coupling at $B_0 = 0$.  Since the background
is exact at $B_0 = 0$, the xPert computation of $L^{(2)}$ involves
no approximation.  $C_0$ is exact.

\paragraph{(v) Complications vanish faster than the signal.}
The graviton--photon coupling in the equations of motion is
\begin{equation}
  g = \bigl[\kappa^2 + \delta_1^2/\alpha + \cdots\bigr] \times B_0
    + \text{corrections}\,.
  \label{eq:bv-coupling}
\end{equation}
The signal (square bracket) is $O(1)$ per unit $B_0$.  Each correction
carries at least one \emph{additional} power of $B_0$ beyond the signal:

\begin{table}[htbp]
\centering
\caption{$B_0$~scaling of signal vs.\ complications.}
\label{tab:bv-scaling}
\begin{tabular}{@{}llll@{}}
\toprule
Contribution & In $g/B_0$ & In $C = P/B_0^2$ & Origin \\
\midrule
\textbf{Signal} ($\kappa^2 + \delta_1^2/\alpha$)
  & $O(1)$ & $O(1)$ & Direct + torsion-mediated coupling \\
Torsion mass correction ($\bar{T} \neq 0$)
  & $O(\delta_1^3 B_0/\alpha^3)$ & $O(B_0)$ & $\alpha_\mathrm{eff} = \alpha + O(\bar{T})$ \\
Dim-6 EFT correction
  & $O(B_0/\Lambda^2)$ & $O(B_0)$ & Suppressed operators \\
Metric backreaction
  & $O(\kappa^2 B_0^2 D^2)$ & $O(B_0^2)$ & Hwang--Noh: $(L/\lambda_B)^2$ \\
\bottomrule
\end{tabular}
\end{table}

At $B_0 = 0.01$ with $\delta_1/\alpha \sim 1$: the leading correction is
$\sim 1\%$ of $C_0$.  At $B_0 = 10^{-5}$: $\sim 10^{-5}$ of $C_0$.

\paragraph{(vi) Linearisation holds throughout.}
\begin{table}[htbp]
\centering
\caption{Perturbation amplitude constraints and their $B_0$~dependence.}
\label{tab:bv-linearisation}
\begin{tabular}{@{}llll@{}}
\toprule
Condition & Requirement & $B_0$~dep. & Satisfied? \\
\midrule
$h \ll 1$ (metric pert.) & $A \ll 1$ & Independent & $\checkmark$ ($A = 0.1$) \\
$P \ll 1$ (conversion) & $\kappa_\mathrm{eff} B_0 D \ll 1$ & $\propto B_0$ & $\checkmark$ (small $B_0$) \\
$|t| \ll \alpha$ (torsion NL) & $\delta_1 B_0 A \ll \alpha^2$ & $\propto B_0$ & $\checkmark$ (small $B_0$) \\
$|a|/|\bar{A}| \ll 1$ (EM pert./bg) & $\kappa_\mathrm{eff} D A / L \ll 1$ & Independent & $\checkmark$ (reasonable $D/L$) \\
\bottomrule
\end{tabular}
\end{table}

Conditions on $h$ and $|a|/|\bar{A}|$ are $B_0$-independent (set by the
initial amplitude $A$ and geometry $D/L$).  Conditions on $P$ and torsion
nonlinearity are proportional to $B_0$ and automatically satisfied in the
$B_0 \to 0$ limit.

The amplification ratio $C_{0,\text{torsion}}/C_{0,\text{EM}}$ is computed
at the same initial amplitude~$A$, so $O(A^2)$ nonlinear corrections cancel
in the ratio.


%% ================================================================
\subsection{Physical $B_0$ Scales in Simulation Units}
\label{sec:bv-physical-scales}

Simulations use natural units with $\kappa \equiv \sqrt{16\pi G} = 1$.
The physically relevant $B_0$ values in these units are:

\begin{table}[htbp]
\centering
\caption{Physical magnetic fields mapped to simulation units ($\kappa = 1$).}
\label{tab:bv-physical-b0}
\begin{tabular}{@{}lrrl@{}}
\toprule
System & $B_0$ (physical) & $B_{0,\text{sim}}$ ($\kappa = 1$) & Regime \\
\midrule
Lab magnet          & $10\,$T ($10^4\,$G)     & $\sim 10^{-18}$ & Far below all corrections \\
Neutron star        & $10^{13}\,$G            & $\sim 10^{-10}$ & Far below all corrections \\
Extreme magnetar    & $10^{15}\,$G            & $\sim 3 \times 10^{-9}$ & Below all corrections \\
\midrule
Simulation (small)  & ---                     & $10^{-5}$       & $10^4 \times$ magnetar \\
Simulation (medium) & ---                     & $10^{-3}$       & $10^6 \times$ magnetar \\
\bottomrule
\end{tabular}
\end{table}

Our simulation values ($B_0 \sim 10^{-3}$ to $10^{-5}$) are
$10^4$--$10^6$ times larger than any physical magnetic field.  This means:
\begin{itemize}
  \item They are ``small'' for backreaction purposes: corrections
    $\sim B_0 \times \delta_1/\alpha \ll 1$.
  \item But orders of magnitude above any real field --- we are solidly
    in the regime where $C = P/B_0^2$ has converged to~$C_0$.
  \item The multi-$B_0$ convergence check (verifying $C$ is constant
    across $B_0$ values) confirms this quantitatively.
\end{itemize}


%% ================================================================
\subsection{Background Torsion: $\bar{T} \neq 0$ for Non-Minimal Theories}
\label{sec:bv-background-torsion}

For the non-minimal coupling $\delta_1 \widetilde{R}_{[\mu\nu]} F^{\mu\nu}$,
the torsion field equation at zeroth order gives:
\begin{equation}
  \alpha\, \bar{T} = -\delta_1 \times f(\bar{F})\,,
  \label{eq:bv-tbar}
\end{equation}
so $\bar{T} \neq 0$ when $B_0 \neq 0$.  Setting $\bar{T} = 0$
(as the current pipeline does) introduces an error.

The error modifies the effective torsion mass:
$\alpha_\mathrm{eff} = \alpha + O(\bar{T}) = \alpha + O(\delta_1 B_0/\alpha)$.
This changes the torsion-mediated coupling from $\delta_1^2/\alpha$ to
$\delta_1^2/\alpha_\mathrm{eff} \approx \delta_1^2/\alpha \times
(1 - \delta_1 B_0/\alpha^2 + \cdots)$.  The correction enters
$g/B_0$ at $O(B_0)$, hence $C = P/B_0^2$ at $O(B_0)$.

\textbf{Resolution:} in the $B_0 \to 0$ limit, $\bar{T} \to 0$
and the background becomes exactly self-consistent.  No need to solve
for~$\bar{T}$.  The conversion coefficient $C_0$ is computed correctly
with $\bar{T} = 0$.

For other non-minimal couplings ($T^2 \times F$, $\nabla T \times F$):
similar sourcing analysis applies.  Each sources $\bar{T}$ at $O(B_0)$,
contributing corrections to $C$ at $O(B_0)$ that vanish in the limit.


%% ================================================================
\subsection{Comparison with Hwang \& Noh}
\label{sec:bv-hwang-noh}

\begin{table}[htbp]
\centering
\caption{Comparison of backreaction arguments.}
\label{tab:bv-comparison}
\begin{tabular}{@{}lll@{}}
\toprule
& Hwang \& Noh (2310.04150) & This work \\
\midrule
Scope & Einstein--Maxwell only & Any PGT+EM theory \\
Regime & Finite $B_0$, bounds error & $B_0 \to 0$ limit, eliminates error \\
Key result & $(L/\lambda_B)^2 \sim P/(2\pi^2)$ & $C_0 = \lim P/B_0^2$ exact \\
Background & Approximate (flat + $B_0$) & Exact (flat, no fields) \\
Torsion & Not considered & Included (non-minimal, algebraic) \\
EFT & Not discussed & Dimension-4 complete, dim-6 suppressed \\
\bottomrule
\end{tabular}
\end{table}

The two arguments are complementary: Hwang--Noh bounds the error at finite
$B_0$ for Einstein--Maxwell; our argument eliminates it entirely in the
$B_0 \to 0$ limit for any PGT+EM theory.


%% ================================================================
\subsection{Verification Protocol}
\label{sec:bv-verification}

For each new theory added to the sweep programme:

\begin{enumerate}
  \item \textbf{Zero-$B_0$ sanity check:} Run at $B_0 = 0$.
    Confirm $P = 0$ identically (no background-independent conversion).
    If $P \neq 0$: the measurement $C = P/B_0^2$ is not valid for this
    theory; investigate the $B_0$-independent channel separately.

  \item \textbf{Multi-$B_0$ convergence:} Run at
    $B_0 \in \{10^{-3}, 10^{-4}, 10^{-5}\}$.
    Compute $C = P/B_0^2$ at each value.  Confirm convergence to $C_0$
    within machine precision (modal solver) or solver tolerance.
    The convergence rate should be $|C(B_0) - C_0| \lesssim
    (\delta_1/\alpha)\, B_0$ (linear in $B_0$).

  \item \textbf{Analytical cross-check} (where available):
    For non-propagating torsion with algebraic constraint:
    $C_0 = (\kappa^2 + \delta_1^2/\alpha)^2 \times D^2/4$.
    Compare numerical $C_0$ against this formula.

  \item \textbf{Linearisation check:} Verify $P \ll 1$ at the chosen
    $B_0$.  If $P > 0.01$, reduce $B_0$ or shorten simulation time.
\end{enumerate}


%% ================================================================
\subsection{Implications for the Sweep Programme}
\label{sec:bv-sweep-implications}

\subsubsection{Primary Metric}

The conversion coefficient $C_0 = \lim_{B_0 \to 0} P/B_0^2$ is the
universal observable for torsion parameter sweeps.  It captures the
intrinsic coupling strength of the theory, independent of the specific
$B_0$ value or background geometry.

\subsubsection{Amplification Factor}

\begin{equation}
  A = \frac{\kappa_\mathrm{eff}}{\kappa}
    = \sqrt{\frac{C_{0,\text{torsion}}}{C_{0,\text{EM}}}}
    = \frac{B_\mathrm{min}(\text{EM})}{B_\mathrm{min}(\text{torsion})}\,.
  \label{eq:bv-amplification}
\end{equation}
The last equality shows that $A$ is identical to the $B$-field lowering
factor: ``torsion reduces the required magnetic field by factor~$A$ for
the same conversion probability.''

\subsubsection{Frequency Scan for Propagating Torsion}

For propagating torsion with mass $\mu = \sqrt{\alpha/b_5}$, the
conversion coefficient $C_0(\omega_0)$ depends on the initial
gravitational-wave frequency~$\omega_0$ relative to~$\mu$.  Scanning
$k_0$ (setting $\omega_0 = k_0$) maps the torsion mass detuning and
distinguishes propagating from non-propagating torsion.

\subsubsection{Simulation Conventions}

Natural units $\hbar = c = 1$, $\kappa = 1$.  Plane-wave graviton IC
in periodic domain with uniform~$B_0$:
\begin{verbatim}
  --ic plane-wave --ic-wavevector k0 --ic-amplitude 0.1 --ic-component h_7
\end{verbatim}
Source~$h_7$ ($h_+$ polarisation), target~$a_2$ ($a_y$ transverse photon),
$k_0 \gg B_0$ (massless dispersion limit), $B_0 \sim 10^{-3}$ to $10^{-5}$.
